\documentclass[book.tex]{subfiles}

\begin{document}
\chapter{Cosmological N-Body Simulations}
\label{chap:n_body}
\noindent\rule{11cm}{0.4pt} \\
\textbf{Prerequisites:} \autoref{chap:molecular_hamiltonian} \\
\textbf{Difficulty Level:} ** \\
\noindent\rule{11cm}{0.4pt}
\vspace{5mm}

N-body simulations provide powerful tools in cosmology to study models of dark matter. These simulations can require evolving billions of particles over very long timescales.

\section{The Particle Mesh Scheme}

This scheme gets the $(N, 3)$ input tensor of the positions of $N$ particles. 

Time evolution is given by the equation

\begin{align}
\frac{dx}{da} &= \frac{1}{a^3 E(a)} v \\
\frac{dv}{da} &= \frac{1}{a^2 E(a)}  F_\theta(x, a) 
\end{align}

\begin{align}
F_\theta(x, a) &= \frac{3 \Omega_m}{2} \nabla \left [ \phi_{PM}(x) * \mathcal{F}^{-1}(1+f_\theta(a, |k|)) \right ]
\end{align}

\section{Convolutional Approximations}

The core idea is to use a convolutional network to approximate the rapid time evolution of the simulation.
\end{document}